What is wealth? A gold ring, a bride-price, a shared life---which of these is the truest form of riches? This paper takes that question seriously and argues that the answer is not economic but ontological: the form of wealth that organises intimate relations determines whether relational subjects can genuinely exist as co-evolving beings, or whether one or both parties are structurally reduced to exchangeable units within a system of accumulation.

The paper develops the concept of \emb{Generative Relational Wealth} (GRW): wealth that is generated through co-evolutionary activity in the relational field, that belongs to the \emph{between} rather than to either individual, and whose most fundamental property is that \emb{use does not deplete but may augment}---because GRW is stored not in any material or symbolic medium but in the permanent structural modifications of the relational attractor landscape. GRW is the self-fuel of the relational field: generated through co-evolution, it becomes the dynamic condition for further co-evolution, constituting what we call the \emph{generative relational cycle}.

Against this concept, the paper maps the major traditions of wealth theory---from mercantilism and classical political economy through Marx's commodity fetishism and Federici's social reproduction theory, from Sen's capabilities approach through Mauss's gift economy and Eglash's generative justice---showing in each case what the tradition illuminates and what it misses about wealth in intimate relations. The paper then examines how traditional wealth forms (bride-price, gold, dowry) embed intimate relations within power-structured logics of accumulation, and how three mechanisms---settlement, commodification, and indebtedness---structurally foreclose genuine relational co-evolution.

A central chapter analyses the \emph{happiness jar}---a thought experiment and practical paradigm in which couples record and process shared moments of happiness---as a case study in GRW accumulation, depletion-resistance, and regeneration. The analysis draws on semiotics (Peirce's index, Benjamin's aura), psychoanalysis (Lacanian drive theory, the objet petit a), and theoretical physics (relational STDP, phase-space reconstruction, the $\Delta$ between symbolic and real registers) to explain why artistic processing of shared memories can reignite a withering relational cycle---and advances the conjecture, clearly marked as speculative, that the operator $O(x) \sim x'$ generates a non-zero relational remainder $\mathcal{R}(O(x), x) \neq 0$ that is the source of GRW's inexhaustibility.

The paper also develops feminist, psychoanalytic, epistemic-justice, and Hegelian-Marxist perspectives on GRW, arguing that: GRW is structurally non-phallic wealth that bypasses the symbolic order's power structures; care labour is GRW's primary production mechanism and relational subjectivity its highest form; the recognition of epistemic asymmetry in intimate relations is a condition for GRW generation; and the dialectical relation between material wealth (necessary condition, secondary aspect) and GRW (primary aspect, generative ground) is irreducible---without symbolic wealth, GRW cannot be sustained; without GRW, symbolic wealth is ontologically impoverished.

This is \pinkref{Paper~XV} of the \emph{Philosophy of Intimacy and the Theory of Justice} series, extending the eudaimonics of \pinkref{Paper~XIV} \citep{paperxiv} and the generativity theory of \pinkref{Paper~IX} \citep{paperix} into the domain of political economy.
